\noindent\textbf{Draft status---author confirmation required.}
The availability statements below describe current sources and local artifacts. Final
author contributions, competing interests, funding, institutional ethics and consent
wording, release accessions and submission metadata remain author-controlled blocking
items; this draft does not represent them as finalized declarations.

\section*{Data Availability}

The public cohorts used in this study retain their source-specific access and reuse terms.
NADT-Prostate is available from The Cancer Imaging Archive
(DOI: 10.7937/TCIA.JHQD-FR46) \cite{nadt2021}; PANDA data are distributed through the
PANDA challenge under its platform terms \cite{panda2022}; TCGA-PRAD histology and
molecular data are available through the Genomic Data Commons, with the referenced
TCGA clinical endpoint fields available through cBioPortal; and PRECISE is available from
Zenodo (DOI: 10.5281/zenodo.20721779) \cite{precise2026}. LEOPARD is an
access-governed challenge resource: its training data and outcome labels remain subject to
the LEOPARD challenge and dataset terms, and its public challenge report is cited here
\cite{leopard2026}. Publication-facing aggregate numerical figure-source tables are supplied
with the public code repository. Patient-level derived analysis artifacts and local provenance
manifests are not redistributed. No separate derived-data archive DOI has yet been assigned;
the final deposition and accession statement must be confirmed before submission.

\section*{Code Availability}

Source code for the audited analyses and manuscript generation is publicly available in
the manuscript development branch at
\url{https://github.com/AiXLab-GNU/vlm-pathology/tree/feat/precise-pni-morphology-rereview}.
Principal auditable analysis entry points are
\begin{itemize}
\item \path{models/build_revision_p0_artifacts.py};
\item \path{models/aggregate_stability_grid.py};
\item \path{models/build_tcga_cdr_pfi_evidence.py};
\item \path{models/build_marker7_survival_paired_analysis.py};
\item \path{models/run_marker7_common_source_sensitivity.py}; and
\item \path{models/build_ar_spop_evidence_closure.py}.
\end{itemize}
The repository also versions the stability runners, manuscript builder, figure renderers,
table generators, validation tests and environment specifications. Manuscript figures are
regenerated from saved CSV source tables by the scripts in \path{paper/figures/}.
Repository-authored software and documentation are released under the MIT License.
Dataset copies, whole-slide images, cached embeddings, pretrained model weights,
access-governed LEOPARD artifacts and patient-level derived outputs are not redistributed;
they must be obtained or reconstructed from their original sources under the applicable
access terms and are not covered by the repository's software license. No versioned archival
DOI was available at the time of manuscript preparation.

\section*{Author Contributions}

No contribution roles were inferred from repository contents. A contribution statement
using the final author list and author-approved roles must replace this factual draft note
before submission.

\section*{Competing Interests}

No competing-interest status was inferred. Each named author's explicit declaration of
interests, including an explicit no-interest statement where applicable, must be confirmed
before submission.

\section*{Ethics and consent statement}

This work was a secondary analysis of deidentified data and involved no new participant
recruitment, intervention or specimen collection. The original participant approvals,
consent procedures and data-use conditions remain those reported by the source publications
and repositories. The corresponding author must confirm the applicable institutional
determination for this secondary analysis and the final approval, waiver and consent wording
before submission; no institutional decision was inferred in this draft.
